\chapter{Osnove računarskih arhitektura i asemblera}
\label{ch:osnove}

\napomena{Teorijsko poglavlje, 8--10 strana. Glavni izvori: \textcite{patterson2020},
\textcite{hennessy2019}, \textcite{harris2022}, \textcite{waterman2019}.}

\section{Arhitektura računarskih sistema}
\label{sec:arhitektura-sistema}
\napomena{Fon Nojmanov model, procesor--memorija--U/I, ciklus dohvati--dekodiraj--izvrši.}

\section{Mašinski jezik i asemblerski jezik}
\label{sec:masinski-asemblerski}
\napomena{Binarni zapis instrukcija; mnemonici, simboličke adrese; odnos 1:1 između
asemblerske i mašinske instrukcije (osim pseudo-instrukcija).}

\section{Instrukcijski skup procesora}
\label{sec:isa}
\napomena{Pojam ISA kao ugovora između hardvera i softvera; vrste instrukcija
(aritmetičke, logičke, memorijske, kontrola toka).}

\section{RISC i CISC arhitekture}
\label{sec:risc-cisc}
\napomena{Istorijat, karakteristike, primjeri (x86 vs.\ ARM/MIPS/RISC-V), load/store princip.}

\section{Pregled RISC-V arhitekture}
\label{sec:risc-v}
\napomena{Nastanak na Berkliju \parencite{asanovic2014,waterman2016}, otvorenost, modularnost
(osnovni skup RV32I/RV64I + ekstenzije M, A, F, D, C), zašto je izabrana za ovaj rad.}

\section{Uloga asemblera u procesu prevođenja programa}
\label{sec:uloga-asemblera}
\napomena{Lanac izvorni kod → kompajler → asembler → objektni fajl → linker → izvršni fajl.
Može dijagram (slika).}
